\section{Sixth Wound: The Saturated Shadow}

The five wounds above were handed to me. Most are the historical objections, raised by people who wanted the mechanical program to succeed, and every one of them is older than I am. This one nobody handed me. I found it while writing the section that closes this chapter, and I have put it here, in the ledger with the others, rather than leaving it where I found it.

Begin with something the shadow push has needed from the start and has so far been allowed to have quietly. Gravity, in this sketch, is a push from a rain that is blocked: a body pushes its neighbour because it removes some of the flux that would otherwise have arrived from that direction. For the push to come out proportional to \emph{mass} --- which is what we observe, and to great precision --- each massion must have only a very small chance of being stopped by any given piece of matter. Matter has to be nearly transparent to the crowd. If it were not --- if a body stopped a large fraction of what reached it --- the far side of the body would sit in the near side's shadow and contribute nothing, and the push would stop tracking the mass and begin tracking the \emph{silhouette}.

Ordinary matter is obliging about this. The Earth is a ball of rock eight thousand miles thick and its gravity still counts every atom, which is a direct measurement, if you like, of how nearly transparent matter is to whatever gravity is made of. Fine. But now ask about a body that is not obliging: something collapsed, something dense enough to trap light. Is it still transparent to the crowd?

If it is not, the arithmetic goes wrong, and it is worth being careful about \emph{how} wrong, because the first version of this section got it wrong in the toy's disfavour and I would rather be accurate than frightening. There is a famous relation in which a black hole's radius grows in direct proportion to its mass, and I reached for it. But that relation describes a \emph{horizon}, and a horizon is not what does any blocking. Blocking is done by matter. And a horizon is a poor stand-in for matter, as its own numbers admit: the average density inside one falls as the inverse square of the mass, so that a hole of ten suns encloses matter at roughly nuclear density, the hole at this galaxy's centre encloses an average density near that of a heavy oil, and a hole of ten thousand million suns encloses an average thinner than the air in this room. Nothing thinner than air is holding anything up. Whatever sits inside a large hole must sit very far inside it, on a strut of its own, at a density set by that strut and not by the horizon.

So take the body seriously and the arithmetic changes. A body held at some characteristic stiffness has a radius growing as the cube root of its mass --- we live in three dimensions, and mass fills volume --- and a silhouette therefore growing as about the two-thirds power. Possibly weaker: real degenerate bodies are perverse about this, and white dwarfs actually \emph{shrink} as they gain mass. So an opaque collapsed body would push as roughly the two-thirds power of its mass, while everything we can weigh insists that the push grows as the mass itself. That is still fatal, and it is fatal in a way that accumulates: the two rules diverge as the cube root of the mass ratio, which is a factor of about seventy-five between a hole of ten suns and the one at the centre of this galaxy, and a factor of about a thousand across the whole observed range of holes. And nature has been unkind about where it set the examination, because collapsed bodies are precisely where we now measure gravity best. The mass of the dark object at our galaxy's centre can be read from the orbits of the stars around it, and from the size of the ring the Event Horizon Telescope photographed, and from nothing disagreeing anywhere; the waveforms of inspiralling pairs match general relativity to many decimal places using one mass throughout. A push that tracked silhouette rather than mass would show up as those numbers refusing to agree with each other. They agree.

Put that way, the wound stops being about black holes at all and becomes a question about \emph{density}, which is an improvement, because a question about density is one the world can be asked. Matter is verified transparent to the crowd wherever gravity has been seen to track mass, and that now includes neutron stars: binary pulsars keep time to a precision that would expose any failure, at densities around four times ten to the seventeenth kilograms in every cubic metre. So transparency holds through nuclear density and out the other side. The question is whether there is any stopping-chance for a massion that leaves matter transparent there and yet turns it opaque at the densities inside a collapsed body --- and whether such a thing could be arranged without looking arranged.

The best reply I have is real but it is not free. A horizon is our layer's \emph{optical} fact: it is defined by what our light cannot climb out of. The crowd is not our light. It belongs to the layer below, it moves very much faster, and it is under no obligation to be stopped by whatever stops a photon. A collapsed body may therefore be black and yet nearly transparent --- opaque to us, and hardly there at all as far as the rain is concerned. Then the push tracks the mass, the orbits come out right, and the wound closes.

But it closes by paying elsewhere, and I want the reader to see the coin leave my hand. The next-to-last section of this chapter will note that black holes carry an entropy proportional to the \emph{area} of the horizon, and will offer the tower an explanation of it. There are two explanations available. The cheap one asks only that our layer be unable to see past the horizon, and it survives this reply intact. The satisfying one --- that the crowd itself is stopped at the surface, so that the shell is genuinely all there is to interact with, the interior having been pressed dense by the shadow of its own shell --- is exactly the opacity this reply has just denied. And the same correction that rescued the arithmetic above charges the satisfying account a second time, which I record because it is my own preferred picture being billed: the entropy in question goes as the area of the \emph{horizon}, which grows as the square of the mass, while a shielding shell at the body's own surface grows only as the two-thirds power. Those are different laws. To have the mechanical account, one would need the body to fill its horizon after all --- and we have just seen what that would require of its density. So the dial has two ends and neither is comfortable. Turn it to opaque: gravity goes wrong by a factor of a thousand across the observed holes, and the area law becomes mechanical --- though, as noted, of the wrong shape. Turn it to transparent: gravity is saved and the area law falls back to bookkeeping. It is the first wound's shape all over again --- a pincer, with the same two jaws closing on the same one knob --- and it is the newest and least settled thing in this book.

What would settle it is a number: the massion's chance of being stopped, per unit of matter it passes through, at densities far beyond anything we can make. Everything above turns on where that number puts the saturation, and I do not know where it puts it. That is the sixth problem of the research program below, and of the six it is the one I would most like to see somebody take away from me.

